%!TEX program = xelatex
\documentclass[11pt,a4paper]{article}
\usepackage[margin=2.2cm]{geometry}
\usepackage{amsmath,amssymb,amsthm,mathtools}
\usepackage{bm}
\usepackage{enumitem}
\usepackage{booktabs}
\usepackage{xcolor}
\usepackage{hyperref}
\usepackage{fancyhdr}
\usepackage{xeCJK}

\hypersetup{
    colorlinks=true,
    linkcolor=blue!60!black,
    citecolor=blue!60!black,
    urlcolor=blue!60!black
}

\pagestyle{fancy}
\fancyhf{}
\rhead{\small Qiuzhen QE Probability \& Statistics --- Fall 2024}
\lhead{\small Official Qualifying Exam}
\rfoot{\small Page \thepage}

\DeclareMathOperator*{\argmax}{arg\,max}
\DeclareMathOperator*{\argmin}{arg\,min}
\DeclareMathOperator{\Var}{Var}
\DeclareMathOperator{\E}{\mathbb{E}}
\DeclareMathOperator{\Pbb}{\mathbb{P}}
\DeclareMathOperator{\Cov}{Cov}

\setlist[itemize]{topsep=3pt,itemsep=2pt,parsep=0pt}
\setlist[enumerate]{topsep=3pt,itemsep=2pt,parsep=0pt}

\newcommand{\examstatement}[1]{%
  \par\addvspace{0.85em}%
  \noindent{\bfseries #1}\par\nobreak\addvspace{0.28em}%
}

% Shared amsthm note/remark/theorem styles
% Shared math extras for qe-review.
% Requires already-loaded: amsmath, amssymb.
%
% Classic pitfall: AMS blackboard-bold fonts have no digit glyphs, so
% \mathbb{1} renders as overlapping garbage. Map the digit 1 to dsfont's
% \mathds{1}; leave \mathbb{R}, \mathbb{P}, etc. on AMS.
%
% Usage (path relative to the main .tex being compiled):
%   notes:  % Shared math extras for qe-review.
% Requires already-loaded: amsmath, amssymb.
%
% Classic pitfall: AMS blackboard-bold fonts have no digit glyphs, so
% \mathbb{1} renders as overlapping garbage. Map the digit 1 to dsfont's
% \mathds{1}; leave \mathbb{R}, \mathbb{P}, etc. on AMS.
%
% Usage (path relative to the main .tex being compiled):
%   notes:  % Shared math extras for qe-review.
% Requires already-loaded: amsmath, amssymb.
%
% Classic pitfall: AMS blackboard-bold fonts have no digit glyphs, so
% \mathbb{1} renders as overlapping garbage. Map the digit 1 to dsfont's
% \mathds{1}; leave \mathbb{R}, \mathbb{P}, etc. on AMS.
%
% Usage (path relative to the main .tex being compiled):
%   notes:  \input{../../common/qe-math-extras.tex}
%   exams:  \input{../../../common/qe-math-extras.tex}

\usepackage{dsfont}
\usepackage{pdftexcmds}

\makeatletter
\let\qe@amsmmathbb\mathbb
\renewcommand{\mathbb}[1]{%
  \ifnum\pdf@strcmp{\unexpanded{#1}}{1}=\z@
    \mathds{1}%
  \else
    \qe@amsmmathbb{#1}%
  \fi
}
\makeatother

% Preferred indicator macros (either style is safe after the remap above).
\newcommand{\bbone}{\mathds{1}}
\newcommand{\ind}{\mathbf{1}}

%   exams:  % Shared math extras for qe-review.
% Requires already-loaded: amsmath, amssymb.
%
% Classic pitfall: AMS blackboard-bold fonts have no digit glyphs, so
% \mathbb{1} renders as overlapping garbage. Map the digit 1 to dsfont's
% \mathds{1}; leave \mathbb{R}, \mathbb{P}, etc. on AMS.
%
% Usage (path relative to the main .tex being compiled):
%   notes:  \input{../../common/qe-math-extras.tex}
%   exams:  \input{../../../common/qe-math-extras.tex}

\usepackage{dsfont}
\usepackage{pdftexcmds}

\makeatletter
\let\qe@amsmmathbb\mathbb
\renewcommand{\mathbb}[1]{%
  \ifnum\pdf@strcmp{\unexpanded{#1}}{1}=\z@
    \mathds{1}%
  \else
    \qe@amsmmathbb{#1}%
  \fi
}
\makeatother

% Preferred indicator macros (either style is safe after the remap above).
\newcommand{\bbone}{\mathds{1}}
\newcommand{\ind}{\mathbf{1}}


\usepackage{dsfont}
\usepackage{pdftexcmds}

\makeatletter
\let\qe@amsmmathbb\mathbb
\renewcommand{\mathbb}[1]{%
  \ifnum\pdf@strcmp{\unexpanded{#1}}{1}=\z@
    \mathds{1}%
  \else
    \qe@amsmmathbb{#1}%
  \fi
}
\makeatother

% Preferred indicator macros (either style is safe after the remap above).
\newcommand{\bbone}{\mathds{1}}
\newcommand{\ind}{\mathbf{1}}

%   exams:  % Shared math extras for qe-review.
% Requires already-loaded: amsmath, amssymb.
%
% Classic pitfall: AMS blackboard-bold fonts have no digit glyphs, so
% \mathbb{1} renders as overlapping garbage. Map the digit 1 to dsfont's
% \mathds{1}; leave \mathbb{R}, \mathbb{P}, etc. on AMS.
%
% Usage (path relative to the main .tex being compiled):
%   notes:  % Shared math extras for qe-review.
% Requires already-loaded: amsmath, amssymb.
%
% Classic pitfall: AMS blackboard-bold fonts have no digit glyphs, so
% \mathbb{1} renders as overlapping garbage. Map the digit 1 to dsfont's
% \mathds{1}; leave \mathbb{R}, \mathbb{P}, etc. on AMS.
%
% Usage (path relative to the main .tex being compiled):
%   notes:  \input{../../common/qe-math-extras.tex}
%   exams:  \input{../../../common/qe-math-extras.tex}

\usepackage{dsfont}
\usepackage{pdftexcmds}

\makeatletter
\let\qe@amsmmathbb\mathbb
\renewcommand{\mathbb}[1]{%
  \ifnum\pdf@strcmp{\unexpanded{#1}}{1}=\z@
    \mathds{1}%
  \else
    \qe@amsmmathbb{#1}%
  \fi
}
\makeatother

% Preferred indicator macros (either style is safe after the remap above).
\newcommand{\bbone}{\mathds{1}}
\newcommand{\ind}{\mathbf{1}}

%   exams:  % Shared math extras for qe-review.
% Requires already-loaded: amsmath, amssymb.
%
% Classic pitfall: AMS blackboard-bold fonts have no digit glyphs, so
% \mathbb{1} renders as overlapping garbage. Map the digit 1 to dsfont's
% \mathds{1}; leave \mathbb{R}, \mathbb{P}, etc. on AMS.
%
% Usage (path relative to the main .tex being compiled):
%   notes:  \input{../../common/qe-math-extras.tex}
%   exams:  \input{../../../common/qe-math-extras.tex}

\usepackage{dsfont}
\usepackage{pdftexcmds}

\makeatletter
\let\qe@amsmmathbb\mathbb
\renewcommand{\mathbb}[1]{%
  \ifnum\pdf@strcmp{\unexpanded{#1}}{1}=\z@
    \mathds{1}%
  \else
    \qe@amsmmathbb{#1}%
  \fi
}
\makeatother

% Preferred indicator macros (either style is safe after the remap above).
\newcommand{\bbone}{\mathds{1}}
\newcommand{\ind}{\mathbf{1}}


\usepackage{dsfont}
\usepackage{pdftexcmds}

\makeatletter
\let\qe@amsmmathbb\mathbb
\renewcommand{\mathbb}[1]{%
  \ifnum\pdf@strcmp{\unexpanded{#1}}{1}=\z@
    \mathds{1}%
  \else
    \qe@amsmmathbb{#1}%
  \fi
}
\makeatother

% Preferred indicator macros (either style is safe after the remap above).
\newcommand{\bbone}{\mathds{1}}
\newcommand{\ind}{\mathbf{1}}


\usepackage{dsfont}
\usepackage{pdftexcmds}

\makeatletter
\let\qe@amsmmathbb\mathbb
\renewcommand{\mathbb}[1]{%
  \ifnum\pdf@strcmp{\unexpanded{#1}}{1}=\z@
    \mathds{1}%
  \else
    \qe@amsmmathbb{#1}%
  \fi
}
\makeatother

% Preferred indicator macros (either style is safe after the remap above).
\newcommand{\bbone}{\mathds{1}}
\newcommand{\ind}{\mathbf{1}}

% Shared amsthm environments for qe-review (minimal).
% Requires already-loaded: amsthm.
% Safe to \input after \usepackage{amsmath,amssymb,amsthm,...}.
%
% Shared numbering uses aliascnt so cleveref keys off the environment
% name (Lemma, Proposition, Exercise, ...) rather than the parent
% counter (Theorem, Definition, Remark). Without aliascnt, \cref{lem:...}
% prints ``Theorem 9.13'' instead of ``Lemma 9.13''.

\usepackage{aliascnt}

\theoremstyle{plain}
\newtheorem{theorem}{Theorem}[section]

\newaliascnt{lemma}{theorem}
\newtheorem{lemma}[lemma]{Lemma}
\aliascntresetthe{lemma}

\newaliascnt{proposition}{theorem}
\newtheorem{proposition}[proposition]{Proposition}
\aliascntresetthe{proposition}

\newaliascnt{corollary}{theorem}
\newtheorem{corollary}[corollary]{Corollary}
\aliascntresetthe{corollary}

\newaliascnt{claim}{theorem}
\newtheorem{claim}[claim]{Claim}
\aliascntresetthe{claim}

\newaliascnt{fact}{theorem}
\newtheorem{fact}[fact]{Fact}
\aliascntresetthe{fact}

\theoremstyle{definition}
\newtheorem{definition}{Definition}[section]

\newaliascnt{exercise}{definition}
\newtheorem{exercise}[exercise]{Exercise}
\aliascntresetthe{exercise}

\newaliascnt{assumption}{definition}
\newtheorem{assumption}[assumption]{Assumption}
\aliascntresetthe{assumption}

\newaliascnt{notation}{definition}
\newtheorem{notation}[notation]{Notation}
\aliascntresetthe{notation}

\newaliascnt{construction}{definition}
\newtheorem{construction}[construction]{Construction}
\aliascntresetthe{construction}

% Example: document-wide numbering (not per section / chapter)
\newtheorem{example}{Example}

\theoremstyle{remark}
\newtheorem{remark}{Remark}[section]
\newtheorem*{remark*}{Remark}

\newaliascnt{heuristic}{remark}
\newtheorem{heuristic}[heuristic]{Heuristic}
\aliascntresetthe{heuristic}

\newaliascnt{convention}{remark}
\newtheorem{convention}[convention]{Convention}
\aliascntresetthe{convention}

\newaliascnt{warning}{remark}
\newtheorem{warning}[warning]{Warning}
\aliascntresetthe{warning}

\newaliascnt{proofsketch}{remark}
\newtheorem{proofsketch}[proofsketch]{Proof sketch}
\aliascntresetthe{proofsketch}

\newaliascnt{checkpoint}{remark}
\newtheorem{checkpoint}[checkpoint]{Checkpoint}
\aliascntresetthe{checkpoint}

% Note: document-wide numbering (not per section)
\newtheorem{note}{Note}
\newtheorem*{note*}{Note}

\newaliascnt{examnote}{note}
\newtheorem{examnote}[examnote]{Exam note}
\aliascntresetthe{examnote}

\newaliascnt{study}{note}
\newtheorem{study}[study]{Study note}
\aliascntresetthe{study}

\newaliascnt{summary}{note}
\newtheorem{summary}[summary]{Summary}
\aliascntresetthe{summary}

\newtheorem{examproblem}{Problem}[section]
\newtheorem*{examproblem*}{Problem}

% Bilingual aliases
\newenvironment{dingli}{\begin{theorem}}{\end{theorem}}
\newenvironment{yinli}{\begin{lemma}}{\end{lemma}}
\newenvironment{mingti}{\begin{proposition}}{\end{proposition}}
\newenvironment{tuilun}{\begin{corollary}}{\end{corollary}}
\newenvironment{dingyi}{\begin{definition}}{\end{definition}}
\newenvironment{li}{\begin{example}}{\end{example}}
\newenvironment{zhu}{\begin{remark}}{\end{remark}}
\newenvironment{biji}{\begin{note}}{\end{note}}
\newenvironment{kaoshi}{\begin{examnote}}{\end{examnote}}

\renewcommand{\qedsymbol}{\ensuremath{\square}}

% PDF sidebar outline + printed TOC for QE exam transcripts.
% Load in the preamble AFTER hyperref and AFTER \newcommand{\examstatement}.
\hypersetup{
  unicode=true,
  bookmarks=true,
  bookmarksopen=true,
  bookmarksopenlevel=3,
  bookmarksnumbered=false,
  bookmarksdepth=3
}
\IfFileExists{bookmark.sty}{%
  \usepackage{bookmark}%
  \bookmarksetup{open,openlevel=3,numbered=false,depth=3}%
}{}
% Unnumbered in print, but still written to the TOC / PDF outline
% down to \subsubsection. Starred headings create neither.
\setcounter{secnumdepth}{-1}
\setcounter{tocdepth}{3}
\makeatletter
\@ifundefined{examstatement}{}{%
  \let\qe@origexamstatement\examstatement
  \renewcommand{\examstatement}[1]{%
    \par\phantomsection
    \addcontentsline{toc}{subsection}{#1}%
    \qe@origexamstatement{#1}%
  }%
}
\makeatother


\title{\textbf{QUALIFYING EXAM: PROBABILITY \& STATISTICS}\\[0.5em]\Large FALL 2024}
\date{}
\author{}

\begin{document}

\maketitle
\vspace{-3em}

\tableofcontents

\section{Instructions}
\begin{itemize}
    \item There are 11 problems in this exam (4 pages). You need to choose 8 of them to solve. If you select more than 8, only the first 8 that you have worked on will be graded. Note that 4 of the problems are worth 15 points each and the rest 10 points each.
    \item You must follow all the rules of exam taking. Misconducts will be subject to proper disciplinary actions by the Center.
    \item You must provide all necessary details for full credits. A final answer with no or little explanation/derivation, even if correct, receives a minimal credit.
    \item \( \mathbb{R} \) denotes the set of real numbers and \( \mathbb{N} = \{1, 2, 3, \dots\} \) denotes the set of positive integers. \( \xrightarrow{(d)} \) and \( \stackrel{(d)}{=} \) mean ``converges in distribution'' and ``equal in distribution'', respectively.
\end{itemize}

\section{Statement of the problems}

\examstatement{Problem 1. [10 pts.]}
Let \( U_1, U_2, \dots \) be independent identically distributed (i.i.d.) random variables uniformly distributed on \( [0, 1] \), and define \( S_n = \sum_{k=1}^n U_k \).
\begin{enumerate}[label=(\alph*)]
    \item Calculate \( \mathbb{E}[(S_n)^4] \).
    \item Determine the distribution of \( S_3 \).
\end{enumerate}

\examstatement{Problem 2. [10 pts.]}
Let \( (Y_n)_{n \ge 1} \) be a sequence of real-valued random variables such that

\[
\sqrt{n}(Y_n - a) \xrightarrow[n \to \infty]{(d)} \mathcal{N}(0, \sigma^2),
\]

where \( \mathcal{N}(0, \sigma^2) \) stands for the normal distribution with \( \sigma \neq 0 \).
\begin{enumerate}[label=(\alph*)]
    \item Let \( g : \mathbb{R} \to \mathbb{R} \) be a function such that it is differentiable at \( a \) and \( g'(a) \neq 0 \). Prove that

    \[
    \sqrt{n}(g(Y_n) - g(a)) \xrightarrow[n \to \infty]{(d)} \mathcal{N}\bigl(0, g'(a)^2 \sigma^2\bigr).
    \]
    \item Fix \( p \in (0, 1) \). For \( n \in \mathbb{N} \), let \( Z_n \stackrel{(d)}{=} \mathrm{Bin}(n, p) \) be a binomial random variable. Prove that

    \[
    \sqrt{n} \left( \ln\left(\frac{Z_n}{n}\right) - \ln p \right) \xrightarrow[n \to \infty]{(d)} \mathcal{N}\left(0, \frac{1-p}{p}\right).
    \]
\end{enumerate}

\examstatement{Problem 3. [10 pts.]}
Let \( X \) be a real-valued random variable. Prove that the following properties are equivalent in the sense that the parameters \( K_i > 0 \) appearing in these properties differ from each other by at most an absolute constant factor.
\begin{enumerate}[label=(\alph*)]
    \item The tails of \( X \) satisfy that

    \[
    \mathbb{P}(|X| \ge t) \le 2 \exp(-t/K_1) \qquad \text{for all } t \ge 0.
    \]
    \item The moments of \( X \) satisfy that

    \[
    \mathbb{E}(|X|^p) \le (K_2 p)^p \qquad \text{for all } p \in \mathbb{N}.
    \]
    \item The moment generating function of \( |X| \) is bounded at some point, i.e.,

    \[
    \mathbb{E} \exp(|X|/K_3) \le 2
    \]

    for some \( K_3 > 0 \).
\end{enumerate}
Hint: You can use Stirling's approximations: \( e(n/e)^n \le n! \le e n (n/e)^n \) for all \( n \in \mathbb{N} \) and

\[
n! = (1+o(1))\sqrt{2\pi n}\,(n/e)^n \qquad \text{for large } n.
\]

\examstatement{Problem 4. [10 pts.]}
Let \( \mathbb{T}_d \) be an infinite \( d \)-regular tree, where every node has degree \( d \). On the other hand, let \( \mathcal{T}_b \) be an infinite \( b \)-ary tree with root \( o \). In other words, \( \mathcal{T}_b \) is an infinite tree where every node has \( b \) children nodes and every non-root node has one parent node. Below is an illustration of a binary tree with \( b = 2 \).

\begin{center}
\small
\( o \)\\[0.7em]
\( 0 \qquad\qquad\qquad\qquad\qquad\qquad 1 \)\\[0.7em]
\( 00 \qquad\qquad 01 \qquad\qquad\qquad 10 \qquad\qquad 11 \)\\[0.7em]
\( 000 \quad 001 \quad 010 \quad 011 \quad 100 \quad 101 \quad 110 \quad 111 \)
\end{center}

Consider the simple random walk \( X_n \) on \( \mathbb{T}_d \) and simple random walk \( Y_n \) on \( \mathcal{T}_b \), where at each step, the walker moves to the neighbor nodes with equal probability.
\begin{enumerate}[label=(\alph*)]
    \item For each \( d \in \mathbb{N} \) with \( d \ge 2 \), determine whether the simple random walk \( X_n \) on \( \mathbb{T}_d \) is recurrent or transient. Prove your claim.
    \item For each \( b \in \mathbb{N} \), determine whether the simple random walk \( Y_n \) on \( \mathcal{T}_b \) is recurrent or transient. Prove your claim. (Hint: Use part (a).)
\end{enumerate}

\examstatement{Problem 5. [15 pts.]}
Let \( X_1, X_2, \dots \) be i.i.d.\ random variables with exponential distribution: \( \mathbb{P}(X_k > x) = e^{-x} \) for \( x \ge 0 \). Define

\[
M_n := \sum_{k=1}^n \frac{X_k}{k}.
\]

\begin{enumerate}[label=(\alph*)]
    \item Prove that \( (M_n - \ln n)_{n \in \mathbb{N}} \) converges to a limit \( Y \) almost surely.
    \item Prove that, for every \( p \in (0, 1) \), \( \bigl( \exp(p M_n)/n^p \bigr)_{n \in \mathbb{N}} \) converges to a limit \( Z \) in \( L^1 \).
\end{enumerate}

\examstatement{Problem 6. [15 pts.]}
Let \( B_t \) be a one-dimensional (1D) standard Brownian motion started from 0.
\begin{enumerate}[label=(\alph*)]
    \item Consider a Brownian motion \( X_t = x + B_t \) started at some \( x > 0 \). For any \( t > 0 \) and \( b > a > 0 \), compute the probability of \( X_t \in [a, b] \) conditioning on that \( X_t \) does not hit zero between 0 and \( t \), i.e.,

    \[
    \mathbb{P}\Bigl( X_t \in [a, b] \Bigm| \min_{0 \le s \le t} X_s > 0 \Bigr).
    \]
    \item A Brownian bridge \( W_t \) on \( [0, 1] \) is a 1D standard Brownian motion \( B_t \) subject to the condition that \( B_1 = 0 \). In other words, \( W_t = (B_t \mid B_1 = 0) \) is a continuous-time Gaussian process whose probability distribution is the conditional probability distribution of \( B_t \) conditioning on \( B_1 = 0 \). A 1D Gaussian free field (GFF) \( h_t \) on \( [0, 1] \) with zero boundary is a continuous-time Gaussian process subject to the zero boundary condition \( h_0 = h_1 = 0 \) and has zero mean \( \mathbb{E} h_t = 0 \), \( t \in [0, 1] \), and covariances

    \[
    \mathbb{E}(h_t h_s) = G(t, s), \qquad t, s \in [0, 1].
    \]

    Here, \( G(t, s) \) is the Green's function of the Laplace operator \( -\Delta \), i.e., \( G(t, s) \) is the unique continuous function such that for any smooth test function \( f \in C_c^\infty(0, 1) \),

    \[
    \int_0^1 G(t, s) \frac{\partial^2}{\partial t^2} f(t)\, dt = -f(s) \qquad \text{and} \qquad G(0, s) = G(1, s) = 0.
    \]

    Prove that the process \( (W_t : t \in [0, 1]) \) has the same distribution as the process \( (h_t : t \in [0, 1]) \) in the sense that for any fixed \( 0 \le t_1 < t_2 < \dots < t_n \le 1 \) and Borel sets \( O_1, O_2, \dots, O_n \),

    \[
    \mathbb{P}(W_{t_1} \in O_1, \dots, W_{t_n} \in O_n) = \mathbb{P}(h_{t_1} \in O_1, \dots, h_{t_n} \in O_n).
    \]

    (Hint: Find the explicit form of the function \( G(t, s) \) and calculate \( \mathbb{E}(W_t W_s) \).)
\end{enumerate}

\examstatement{Problem 7. [10 pts.]}
Let \( X_1, \dots, X_n \) be an iid sample from \( N(\mu, 1) \) with \( \mu \) unknown. Unfortunately, one forgets to record \( X_1, \dots, X_n \) in a study and only records \( Y = (Y_1, \dots, Y_n) \) where \( Y_i = I(X_i < 0) \) and \( I(\cdot) \) is the indicator function.
\begin{enumerate}[label=(\alph*)]
    \item Derive the MLE of \( \mu \) based on the observed data \( Y \).
    \item Construct a size \( \alpha \) uniformly most powerful (UMP) test for testing \( H_0 : \mu \le \mu_0 \) versus \( H_1 : \mu > \mu_0 \) based on the observed data \( Y \).
    \item Describe how to construct a \( (1 - \alpha) \) confidence interval for \( \mu \) based on the observed data \( Y \).
\end{enumerate}

\examstatement{Problem 8. [10 pts.]}
Consider the following linear model \( Y_i = \mathbf{z}_i^T \beta + \epsilon_i \), \( i = 1, \dots, n \). \( \mathbf{z}_1, \dots, \mathbf{z}_n \in \mathbb{R}^d \) are fixed and given, and \( \beta \in \mathbb{R}^d \) is unknown. \( \epsilon_i \)'s are random variables satisfying the Gauss-Markov assumptions that \( \mathbb{E}[\epsilon_i] = 0 \), \( \mathrm{Var}[\epsilon_i] = \sigma^2 \) and \( \mathrm{Cov}(\epsilon_i, \epsilon_j) = 0 \), \( \forall i \neq j \). Let \( Y = (Y_1, \dots, Y_n)^T \), and

\[
Z =
\begin{pmatrix}
\mathbf{z}_1^T \\
\mathbf{z}_2^T \\
\vdots \\
\mathbf{z}_n^T
\end{pmatrix}
\]

be the \( n \) by \( d \) design matrix.
\begin{enumerate}[label=(\alph*)]
    \item Let \( \hat{\beta} \) be the least squares estimate of \( \beta \) which is given by \( \hat{\beta} = (Z^T Z)^{-1} Z^T Y \). Let \( \theta = \mathbf{b}^T \beta \) where \( \mathbf{b} \in \mathbb{R}^d \) is a known vector. Write down the mean and the variance of \( \hat{\theta} \) where \( \hat{\theta} = \mathbf{b}^T \hat{\beta} \). Further, prove that under the Gauss-Markov assumptions, the estimator \( \hat{\theta} \) has the smallest variance among all linear unbiased estimator of \( \theta \). Here linear unbiased estimator we mean estimator in the form of \( \mathbf{c}^T Y \) and is unbiased for \( \theta \).
    \item Further assume that \( (\epsilon_1, \dots, \epsilon_n) \) are iid from \( N(0, \sigma^2) \) with \( \sigma^2 \) known. Derive the information matrix \( I(\beta) \).
\end{enumerate}

\examstatement{Problem 9. [10 pts.]}
Suppose \( X_1, \dots, X_n \) are IID from the uniform distribution on \( [0, \theta] \) for some unknown \( \theta > 0 \). Fix \( t \in (0, \theta) \). Consider two estimators of \( \mathbb{P}(X_1 \le t) \): \( F_n(t) = \frac{1}{n} \sum_{i=1}^n I_{\{X_i \le t\}} \) and \( T_n(t) = t/(2\overline{X}) \), where \( \overline{X} \) is the sample mean.
\begin{enumerate}[label=(\alph*)]
    \item Find the asymptotic distributions of the two estimators.
    \item For what value of \( t \) will the first estimator have a smaller asymptotic variance than the second estimator?
    \item Let \( \theta = 1 \). For the \( F_n(t) \) defined above, find the asymptotic distribution of \( n F_n(n^{-1/2}) - \sqrt{n} \).
\end{enumerate}

\examstatement{Problem 10. [15 pts.]}
Let \( X_1, \dots, X_n \) (\( n \ge 2 \)) be iid from \( N(\mu, \sigma^2) \) distribution with \( \mu \ge 0 \) and \( \sigma > 0 \) being the unknown parameters. Let \( \overline{X} \) and \( S^2 \) be the sample mean and sample variance, respectively. Recall \( \chi^2_k \) has probability density function

\[
\frac{1}{2^{k/2} \Gamma(k/2)} x^{k/2 - 1} e^{-x/2}, \qquad x \ge 0.
\]

\begin{enumerate}[label=(\alph*)]
    \item Show \( \overline{X} \) and \( S^2 \) are independent.
    \item Find UMVUE of \( \mu/\sigma \) if it exists.
    \item Is \( \overline{X} \) admissible for estimating \( \mu \) under the square error loss? Prove your assertion.
\end{enumerate}

\examstatement{Problem 11. [15 pts.]}
Let \( X_1, \dots, X_n \) be an iid sample from Uniform\( [\theta, \theta + |\theta|] \) where \( \theta \neq 0 \).
\begin{enumerate}[label=(\alph*)]
    \item Derive the method of moments estimator of \( \theta \).
    \item Derive the MLE of \( \theta \), \( \hat{\theta} \).
    \item Is \( \hat{\theta} \) a consistent estimator of \( \theta \)? Please explain your answer.
\end{enumerate}

\noindent\rule{\linewidth}{0.4pt}

\section{Statement and solutions}

\subsection{Problem 1. [10 pts.]}
Let \( U_1, U_2, \dots \) be independent identically distributed (i.i.d.) random variables uniformly distributed on \( [0, 1] \), and define \( S_n = \sum_{k=1}^n U_k \).
\begin{enumerate}[label=(\alph*)]
    \item Calculate \( \mathbb{E}[(S_n)^4] \).
    \item Determine the distribution of \( S_3 \).
\end{enumerate}

\subsection{Problem 2. [10 pts.]}
Let \( (Y_n)_{n \ge 1} \) be a sequence of real-valued random variables such that

\[
\sqrt{n}(Y_n - a) \xrightarrow[n \to \infty]{(d)} \mathcal{N}(0, \sigma^2),
\]

where \( \mathcal{N}(0, \sigma^2) \) stands for the normal distribution with \( \sigma \neq 0 \).
\begin{enumerate}[label=(\alph*)]
    \item Let \( g : \mathbb{R} \to \mathbb{R} \) be a function such that it is differentiable at \( a \) and \( g'(a) \neq 0 \). Prove that

    \[
    \sqrt{n}(g(Y_n) - g(a)) \xrightarrow[n \to \infty]{(d)} \mathcal{N}\bigl(0, g'(a)^2 \sigma^2\bigr).
    \]
    \item Fix \( p \in (0, 1) \). For \( n \in \mathbb{N} \), let \( Z_n \stackrel{(d)}{=} \mathrm{Bin}(n, p) \) be a binomial random variable. Prove that

    \[
    \sqrt{n} \left( \ln\left(\frac{Z_n}{n}\right) - \ln p \right) \xrightarrow[n \to \infty]{(d)} \mathcal{N}\left(0, \frac{1-p}{p}\right).
    \]
\end{enumerate}

\subsection{Problem 3. [10 pts.]}
Let \( X \) be a real-valued random variable. Prove that the following properties are equivalent in the sense that the parameters \( K_i > 0 \) appearing in these properties differ from each other by at most an absolute constant factor.
\begin{enumerate}[label=(\alph*)]
    \item The tails of \( X \) satisfy that

    \[
    \mathbb{P}(|X| \ge t) \le 2 \exp(-t/K_1) \qquad \text{for all } t \ge 0.
    \]
    \item The moments of \( X \) satisfy that

    \[
    \mathbb{E}(|X|^p) \le (K_2 p)^p \qquad \text{for all } p \in \mathbb{N}.
    \]
    \item The moment generating function of \( |X| \) is bounded at some point, i.e.,

    \[
    \mathbb{E} \exp(|X|/K_3) \le 2
    \]

    for some \( K_3 > 0 \).
\end{enumerate}
Hint: You can use Stirling's approximations: \( e(n/e)^n \le n! \le e n (n/e)^n \) for all \( n \in \mathbb{N} \) and

\[
n! = (1+o(1))\sqrt{2\pi n}\,(n/e)^n \qquad \text{for large } n.
\]

\subsection{Problem 4. [10 pts.]}
Let \( \mathbb{T}_d \) be an infinite \( d \)-regular tree, where every node has degree \( d \). On the other hand, let \( \mathcal{T}_b \) be an infinite \( b \)-ary tree with root \( o \). In other words, \( \mathcal{T}_b \) is an infinite tree where every node has \( b \) children nodes and every non-root node has one parent node. Below is an illustration of a binary tree with \( b = 2 \).

\begin{center}
\small
\( o \)\\[0.7em]
\( 0 \qquad\qquad\qquad\qquad\qquad\qquad 1 \)\\[0.7em]
\( 00 \qquad\qquad 01 \qquad\qquad\qquad 10 \qquad\qquad 11 \)\\[0.7em]
\( 000 \quad 001 \quad 010 \quad 011 \quad 100 \quad 101 \quad 110 \quad 111 \)
\end{center}

Consider the simple random walk \( X_n \) on \( \mathbb{T}_d \) and simple random walk \( Y_n \) on \( \mathcal{T}_b \), where at each step, the walker moves to the neighbor nodes with equal probability.
\begin{enumerate}[label=(\alph*)]
    \item For each \( d \in \mathbb{N} \) with \( d \ge 2 \), determine whether the simple random walk \( X_n \) on \( \mathbb{T}_d \) is recurrent or transient. Prove your claim.
    \item For each \( b \in \mathbb{N} \), determine whether the simple random walk \( Y_n \) on \( \mathcal{T}_b \) is recurrent or transient. Prove your claim. (Hint: Use part (a).)
\end{enumerate}

\subsection{Problem 5. [15 pts.]}
Let \( X_1, X_2, \dots \) be i.i.d.\ random variables with exponential distribution: \( \mathbb{P}(X_k > x) = e^{-x} \) for \( x \ge 0 \). Define

\[
M_n := \sum_{k=1}^n \frac{X_k}{k}.
\]

\begin{enumerate}[label=(\alph*)]
    \item Prove that \( (M_n - \ln n)_{n \in \mathbb{N}} \) converges to a limit \( Y \) almost surely.
    \item Prove that, for every \( p \in (0, 1) \), \( \bigl( \exp(p M_n)/n^p \bigr)_{n \in \mathbb{N}} \) converges to a limit \( Z \) in \( L^1 \).
\end{enumerate}

\subsection{Problem 6. [15 pts.]}
Let \( B_t \) be a one-dimensional (1D) standard Brownian motion started from 0.
\begin{enumerate}[label=(\alph*)]
    \item Consider a Brownian motion \( X_t = x + B_t \) started at some \( x > 0 \). For any \( t > 0 \) and \( b > a > 0 \), compute the probability of \( X_t \in [a, b] \) conditioning on that \( X_t \) does not hit zero between 0 and \( t \), i.e.,

    \[
    \mathbb{P}\Bigl( X_t \in [a, b] \Bigm| \min_{0 \le s \le t} X_s > 0 \Bigr).
    \]
    \item A Brownian bridge \( W_t \) on \( [0, 1] \) is a 1D standard Brownian motion \( B_t \) subject to the condition that \( B_1 = 0 \). In other words, \( W_t = (B_t \mid B_1 = 0) \) is a continuous-time Gaussian process whose probability distribution is the conditional probability distribution of \( B_t \) conditioning on \( B_1 = 0 \). A 1D Gaussian free field (GFF) \( h_t \) on \( [0, 1] \) with zero boundary is a continuous-time Gaussian process subject to the zero boundary condition \( h_0 = h_1 = 0 \) and has zero mean \( \mathbb{E} h_t = 0 \), \( t \in [0, 1] \), and covariances

    \[
    \mathbb{E}(h_t h_s) = G(t, s), \qquad t, s \in [0, 1].
    \]

    Here, \( G(t, s) \) is the Green's function of the Laplace operator \( -\Delta \), i.e., \( G(t, s) \) is the unique continuous function such that for any smooth test function \( f \in C_c^\infty(0, 1) \),

    \[
    \int_0^1 G(t, s) \frac{\partial^2}{\partial t^2} f(t)\, dt = -f(s) \qquad \text{and} \qquad G(0, s) = G(1, s) = 0.
    \]

    Prove that the process \( (W_t : t \in [0, 1]) \) has the same distribution as the process \( (h_t : t \in [0, 1]) \) in the sense that for any fixed \( 0 \le t_1 < t_2 < \dots < t_n \le 1 \) and Borel sets \( O_1, O_2, \dots, O_n \),

    \[
    \mathbb{P}(W_{t_1} \in O_1, \dots, W_{t_n} \in O_n) = \mathbb{P}(h_{t_1} \in O_1, \dots, h_{t_n} \in O_n).
    \]

    (Hint: Find the explicit form of the function \( G(t, s) \) and calculate \( \mathbb{E}(W_t W_s) \).)
\end{enumerate}

\subsection{Problem 7. [10 pts.]}
Let \( X_1, \dots, X_n \) be an iid sample from \( N(\mu, 1) \) with \( \mu \) unknown. Unfortunately, one forgets to record \( X_1, \dots, X_n \) in a study and only records \( Y = (Y_1, \dots, Y_n) \) where \( Y_i = I(X_i < 0) \) and \( I(\cdot) \) is the indicator function.
\begin{enumerate}[label=(\alph*)]
    \item Derive the MLE of \( \mu \) based on the observed data \( Y \).
    \item Construct a size \( \alpha \) uniformly most powerful (UMP) test for testing \( H_0 : \mu \le \mu_0 \) versus \( H_1 : \mu > \mu_0 \) based on the observed data \( Y \).
    \item Describe how to construct a \( (1 - \alpha) \) confidence interval for \( \mu \) based on the observed data \( Y \).
\end{enumerate}

\subsection{Problem 8. [10 pts.]}
Consider the following linear model \( Y_i = \mathbf{z}_i^T \beta + \epsilon_i \), \( i = 1, \dots, n \). \( \mathbf{z}_1, \dots, \mathbf{z}_n \in \mathbb{R}^d \) are fixed and given, and \( \beta \in \mathbb{R}^d \) is unknown. \( \epsilon_i \)'s are random variables satisfying the Gauss-Markov assumptions that \( \mathbb{E}[\epsilon_i] = 0 \), \( \mathrm{Var}[\epsilon_i] = \sigma^2 \) and \( \mathrm{Cov}(\epsilon_i, \epsilon_j) = 0 \), \( \forall i \neq j \). Let \( Y = (Y_1, \dots, Y_n)^T \), and

\[
Z =
\begin{pmatrix}
\mathbf{z}_1^T \\
\mathbf{z}_2^T \\
\vdots \\
\mathbf{z}_n^T
\end{pmatrix}
\]

be the \( n \) by \( d \) design matrix.
\begin{enumerate}[label=(\alph*)]
    \item Let \( \hat{\beta} \) be the least squares estimate of \( \beta \) which is given by \( \hat{\beta} = (Z^T Z)^{-1} Z^T Y \). Let \( \theta = \mathbf{b}^T \beta \) where \( \mathbf{b} \in \mathbb{R}^d \) is a known vector. Write down the mean and the variance of \( \hat{\theta} \) where \( \hat{\theta} = \mathbf{b}^T \hat{\beta} \). Further, prove that under the Gauss-Markov assumptions, the estimator \( \hat{\theta} \) has the smallest variance among all linear unbiased estimator of \( \theta \). Here linear unbiased estimator we mean estimator in the form of \( \mathbf{c}^T Y \) and is unbiased for \( \theta \).
    \item Further assume that \( (\epsilon_1, \dots, \epsilon_n) \) are iid from \( N(0, \sigma^2) \) with \( \sigma^2 \) known. Derive the information matrix \( I(\beta) \).
\end{enumerate}

\subsection{Problem 9. [10 pts.]}
Suppose \( X_1, \dots, X_n \) are IID from the uniform distribution on \( [0, \theta] \) for some unknown \( \theta > 0 \). Fix \( t \in (0, \theta) \). Consider two estimators of \( \mathbb{P}(X_1 \le t) \): \( F_n(t) = \frac{1}{n} \sum_{i=1}^n I_{\{X_i \le t\}} \) and \( T_n(t) = t/(2\overline{X}) \), where \( \overline{X} \) is the sample mean.
\begin{enumerate}[label=(\alph*)]
    \item Find the asymptotic distributions of the two estimators.
    \item For what value of \( t \) will the first estimator have a smaller asymptotic variance than the second estimator?
    \item Let \( \theta = 1 \). For the \( F_n(t) \) defined above, find the asymptotic distribution of \( n F_n(n^{-1/2}) - \sqrt{n} \).
\end{enumerate}

\subsection{Problem 10. [15 pts.]}
Let \( X_1, \dots, X_n \) (\( n \ge 2 \)) be iid from \( N(\mu, \sigma^2) \) distribution with \( \mu \ge 0 \) and \( \sigma > 0 \) being the unknown parameters. Let \( \overline{X} \) and \( S^2 \) be the sample mean and sample variance, respectively. Recall \( \chi^2_k \) has probability density function

\[
\frac{1}{2^{k/2} \Gamma(k/2)} x^{k/2 - 1} e^{-x/2}, \qquad x \ge 0.
\]

\begin{enumerate}[label=(\alph*)]
    \item Show \( \overline{X} \) and \( S^2 \) are independent.
    \item Find UMVUE of \( \mu/\sigma \) if it exists.
    \item Is \( \overline{X} \) admissible for estimating \( \mu \) under the square error loss? Prove your assertion.
\end{enumerate}

\subsection{Problem 11. [15 pts.]}
Let \( X_1, \dots, X_n \) be an iid sample from Uniform\( [\theta, \theta + |\theta|] \) where \( \theta \neq 0 \).
\begin{enumerate}[label=(\alph*)]
    \item Derive the method of moments estimator of \( \theta \).
    \item Derive the MLE of \( \theta \), \( \hat{\theta} \).
    \item Is \( \hat{\theta} \) a consistent estimator of \( \theta \)? Please explain your answer.
\end{enumerate}
\end{document}
